% !TeX root = main.tex

\phantomsection
\section*{Résumé}
\addcontentsline{toc}{section}{Résumé}

Les cabinets de géomètre-expert disposent, au fil des décennies, d'un fonds d'archives considérable composé de plans de bornage, procès-verbaux et documents d'arpentage sous des formats très hétérogènes, tous à caractère foncier. Une grande partie de ce patrimoine, produit avant la généralisation des outils numériques, demeure non indexée et non versée sur la plateforme de référence de la profession, Géofoncier. Ce travail vise à concevoir et à déployer un outil automatisé pour traiter ces archives au sein du cabinet GEO-SIAPP, en Ardèche, en vue de leur versement sur la plateforme nationale Géofoncier.

L'outil développé repose sur une chaîne de traitement
articulée autour de quatre étapes : la classification automatique du document entrant, l'extraction des informations clés (commune, date, numéro de document d'arpentage), la fiabilisation des données par croisement avec des référentiels officiels, et la préparation au versement sur Géofoncier. Pour faire face à la grande variabilité des documents et plans cadastraux récents, registres manuscrits anciens, images numérisées, l'architecture mobilise des techniques issues de la vision par ordinateur, de la reconnaissance de texte manuscrit et de l'intelligence artificielle.

Une interface de validation permet aux opérateurs de contrôler et corriger les résultats de la machine avant tout versement définitif, assurant ainsi la qualité de la donnée foncière produite et le respect des obligations déontologiques de la profession.

\vspace{1.5em}
\noindent\textbf{Mots-clés :} archivage numérique, reconnaissance de texte manuscrit (HTR), traitement automatisé, Géofoncier, intelligence artificielle, géomètre-expert, vision par ordinateur.

{\let\thefootnote\relax\footnotetext{\footnotesize Ce rapport a fait l'objet d'une relecture et d'une correction orthographique et grammaticale assistées par intelligence artificielle.}}

\clearpage

\section*{Abstract}
\addcontentsline{toc}{section}{Abstract}

Land surveying firms build up, over decades, a large archive of boundary delimitation plans, survey reports and cadastral documents in highly heterogeneous formats. A significant part of this documentary heritage, produced before digital tools became widely used, remains unindexed and has not been uploaded to Geofoncier, the profession's national platform. This work aims to design and deploy an automated archive processing tool within the GEO-SIAPP firm in Ardèche, France, in order to feed this data into the national Geofoncier platform.

The developed tool relies on a modular processing chain built around four steps: automatic classification of incoming documents, extraction of key information (municipality, date, cadastral survey number), data validation by cross-referencing official databases, and preparation for submission to Geofoncier. To handle the high variability of the documents, recent cadastral plans (old handwritten registers, degraded scanned images) the architecture combines techniques from computer vision, handwritten text recognition and generative artificial intelligence.

A validation interface allows operators to review and correct the machine's results before the final submission, guaranteeing the quality of the land data produced and compliance with the profession's ethical obligations.

\vspace{1.5em}
\noindent\textbf{Keywords:} digital archiving, handwritten text recognition (HTR), automated processing, Geofoncier, artificial intelligence, land surveyor, computer vision.
